%%%%%%%%%%%%%%%%%%%%%%%%%%%%%%%%%%%%%%%%%%%%%%%%%%%%%%%%%%%%%%
%%%%%%%%%%%%%%%%%% MA-0756 Ecuaciones diferenciales parciales %%%%%%%%%%%%%%%%%%%%%%%%%%%%%%
%%%%%%%%%%%%%%%%%%%%%%%%%%%%%%%%%%%%%%%%%%%%%%%%%%%%%%%%%%%%%%
\newpage
\subsection*{MA-0756 Ecuaciones diferenciales parciales}
\addcontentsline{toc}{subsection}{MA-0756 Ecuaciones diferenciales parciales}

\hypertarget{MA0756}{}
\begin{refsection}
\begin{table}[H]
\centering{
\begin{tabular}{|ll|}
\hline
\textbf{Año:} IV  & \textbf{Requisitos:} MA-0788 Análisis Real II; MA-0515 Ecuaciones Diferenciales \\ 
\textbf{Ciclo:} Optativa  & \textbf{Correquisitos:} Ninguno \\ 
\textbf{Tipo de curso:} Teórico & \textbf{Horas presenciales por semana:} 5 \\ 
\textbf{Créditos:} 5 & \textbf{Horas de trabajo independiente por semana:} 10 \\
\textbf{Virtualidad:} P  &  \\ \hline
\end{tabular}
}
\end{table}

\subsubsection*{Descripción del curso}

Este curso optativo presenta una introducción rigurosa a la teoría de ecuaciones diferenciales parciales desde la perspectiva del análisis funcional y de los espacios de Sobolev. Parte de un repaso concentrado de las ecuaciones prototipo clásicas (Laplace, calor y onda) y avanza hacia formulaciones débiles, existencia de soluciones y resultados de evolución lineal.

El programa sigue, en lo esencial, la estructura de un curso de posgrado introductorio (como el de \emph{Partial Differential Equations} de L.~C.~Evans), adaptado al nivel de la licenciatura avanzada: se enfatizan demostraciones, estimaciones de energía y el uso de herramientas como el teorema de Lax--Milgram y aproximaciones de Galerkin. Se supone dominio de medida e integración, espacios $L^p$ y cálculo en varias variables.

\subsubsection*{Objetivos}

\textbf{General:} Desarrollar competencias teóricas para formular, analizar y resolver problemas de ecuaciones diferenciales parciales mediante métodos clásicos y formulaciones débiles.

\textbf{Específicos:} Al finalizar el curso se espera que la persona estudiante sea capaz de:

\begin{enumerate}
\item Recordar y aplicar soluciones fundamentales y métodos de energía para las ecuaciones de Laplace, calor y onda.
\item Trabajar con derivadas débiles y espacios de Sobolev $W^{k,p}$, incluyendo teoremas de traza e imersiones.
\item Demostrar existencia (y unicidad, cuando corresponda) de soluciones débiles para operadores elípticos de segundo orden mediante Lax--Milgram y la alternativa de Fredholm.
\item Describir ideas básicas de regularidad interior en el contexto elíptico.
\item Construir soluciones aproximadas de ecuaciones de evolución lineales (parabólicas o hiperbólicas) por Galerkin y estimaciones de energía.
\item Comunicar argumentos analíticos con rigor, de forma oral y escrita.
\end{enumerate}

\subsubsection*{Contenidos}

\begin{enumerate}

\item \textbf{Ecuaciones clásicas fundamentales (3 semanas; Evans, cap.~2):}
\begin{enumerate}
    \item Repaso intensivo de soluciones fundamentales y métodos de energía para las ecuaciones de Laplace, calor y onda.
\end{enumerate}

\item \textbf{Espacios de Sobolev (4 semanas; Evans, cap.~5):}
\begin{enumerate}
    \item Derivadas débiles y espacios $W^{k,p}$.
    \item Aproximaciones por funciones suaves y teorema de la traza.
    \item Teoremas de imersión de Sobolev y de Morrey.
\end{enumerate}

\item \textbf{Ecuaciones elípticas de segundo orden (4 semanas; Evans, cap.~6):}
\begin{enumerate}
    \item Formulación débil y existencia mediante el teorema de Lax--Milgram.
    \item Alternativa de Fredholm para problemas elípticos.
    \item Introducción breve a regularidad interior.
\end{enumerate}

\item \textbf{Ecuaciones de evolución lineales (3 semanas; Evans, cap.~7):}
\begin{enumerate}
    \item Ecuaciones parabólicas (calor) o hiperbólicas (onda): aproximaciones de Galerkin.
    \item Estimaciones de energía y existencia de soluciones débiles.
\end{enumerate}

\end{enumerate}

\subsubsection*{Metodología}

\noindent \textbf{Lineamientos metodológicos:} La persona docente a cargo de este curso debe respetar las siguientes pautas:

\begin{enumerate}
    \item Combinar \textbf{clases magistrales} con \textbf{sesiones de ejercicios} en las que se discutan demostraciones y problemas tipo.
    \item Exigir trabajo independiente sustantivo para asimilar la teoría de Sobolev y las técnicas de existencia.
    \item Promover la lectura guiada de un texto de referencia (Evans y, en Sobolev, material complementario de Brezis).
    \item Evaluar mediante tareas periódicas, exámenes parciales y, cuando corresponda, un trabajo escrito sobre un tema del programa.
    \item Cuando las autoridades sanitarias o institucionales impongan restricciones, adaptar las actividades según normativa vigente.
\end{enumerate}

\textbf{Sugerencias metodológicas:} Adicionalmente, se tienen las siguientes recomendaciones para la persona docente a cargo de este curso:

\begin{enumerate}
    \item Explicitar las conexiones entre el repaso del capítulo~2 y los capítulos posteriores (formulación débil y evolución).
    \item Señalar qué partes del capítulo~7 se desarrollarán con énfasis parabólico, hiperbólico o ambos, según la oferta del ciclo.
    \item Alternar demostraciones completas con bosquejos que ilustren la estrategia de las pruebas más técnicas.
\end{enumerate}

\nocite{Eva97}
\nocite{Bre11}

\printcursosbibliography

\end{refsection}
